\section{Introduction}

\subsection{Contributions}

This paper makes the following contributions:

The
methodology addresses a gap identified across three mature but previously
separate research traditions: technology roadmapping (TRM),
capability-based planning (CBP), and cybersecurity capability maturity
assessment. To the best of our knowledge, no existing framework integrates
all four elements---(a)~architecture-anchored capability derivation,
(b)~multi-criteria prioritisation, (c)~multi-dimensional profiling across
technical, regulatory, market, and sovereignty axes, and (d)~phased
progression with quantitative maturity metrics---in a unified, traceable
pipeline tailored to the Cloud--Edge--IoT cybersecurity domain.

Canonical TRM frameworks~\cite{phaal2004,garcia1997} provide the phased
planning structure but do not natively encode regulatory, sovereignty, or
market dimensions. Capability-based planning~\cite{davis2002} provides the
outcome-oriented capability concept and DOTMLPF-P decomposition logic but
has not been systematically applied to civilian cybersecurity roadmapping.
EU cybersecurity roadmap precedents---CyberROAD~\cite{ariu2016},
CONCORDIA~\cite{concordia}, CyberSec4Europe~\cite{cybersec4europe_d45}---provide
multi-dimensional framing and phased recommendations but lack quantitative
per-capability maturity metrics and architecture-anchored capability
derivation. The present methodology synthesises these traditions into a
single reproducible instrument.

\begin{itemize}
    \item A structured pipeline methodology for cybersecurity capability prioritization in the Cloud--Edge--IoT continuum.
    \item A 4-dimensional filtering logic linking regulatory, architectural, sovereignty, and market-adoption considerations.
    \item A hybrid 7-dimensional profiling framework for actionable R\&I prioritization.
    \item A progression matrix model capturing operational deployment and strategic sovereignty evolution across time phases.
    \item A traceable transformation from technical capabilities to policy-aligned strategic recommendations.
\end{itemize}
\newpage
% =========================================================================
% CHAPTER 1: INSTRUCTIONS
% =========================================================================

\noindent
This manuscript and roadmap defines the architectural guardrails required to secure the Cloud--Edge--IoT Continuum. It moves beyond traditional ``perimeter security'' toward a model of cyber-resilience---where systems are engineered to proactively anticipate, withstand, recover from, and adapt to attacks while maintaining strict digital sovereignty.
\textbf{Approach:} To provide actionable, risk-aware Research \& Innovation (R\&I) priorities, the strategy is organized into a two-part framework:
 

\subsection*{Purpose of This Document}
This roadmap provides a structured methodology for European policymakers, research organizations, and industry stakeholders to:
\begin{itemize}
    \item Identify and prioritize cybersecurity R\&I investments for the Cloud-Edge-IoT continuum
\item Track capability maturity using objective Operational Deployment Phase (ODP) and Strategic Sovereignty Level (SSL) metrics.
    \item Align Horizon Europe funding with strategic autonomy goals and regulatory compliance requirements
    \item Ensure coordinated implementation across Member States and European initiatives (ECCC, Digital Europe Programme, EU Chips Act)
\end{itemize}

\subsection*{Target Audiences}
\begin{description}
    \item[European Commission \& Funding Bodies:] Use the 4-Filter Prioritization Framework to allocate Horizon Europe, Digital Europe Programme, and ECCC funding to capabilities with the highest strategic impact.
    
    \item[National Cybersecurity Agencies:] Leverage the progression matrices to coordinate national R\&I investments with EU-wide timelines and ensure interoperability across Member States.
    
    \item[Research Consortia \& Industry:] Utilize the 7-Dimensional Analysis template to structure project proposals, clearly articulating baseline, target, drivers, risks, and alignment with European strategic goals.
    
    \item[Standardization Bodies (CEN/CENELEC, ETSI):] Reference the architectural framework and MIMs to develop harmonized standards for secure, interoperable computing continuum infrastructure.
\end{description}

\subsection*{How to Navigate This Document}
\begin{enumerate}
    \item \textbf{Start with Section 1-3:} Understand the threat landscape, regulatory drivers, and strategic context for securing the European computing continuum.
    
    \item \textbf{Review Section 4:} Familiarize yourself with the 6 Security Layers and 4 Cross-Layers comprising the architectural framework.
    
    \item \textbf{Apply Section 5:} Use the 4-Filter Prioritization Framework to assess which capabilities are critical for your specific context (national strategy, industrial sector, research focus).
    
    \item \textbf{Study Section 6-7:} Understand the 7-Dimensional Analysis methodology and ODP/SSL progression metrics before designing projects or policies.
    
    \item \textbf{Reference Section 9 (Priorities 1-3):} Use the detailed capability profiles as templates for structuring additional priorities relevant to your domain.
    
    \item \textbf{Implement Section 10:} Actionable strategic recommendations for Phase 1 (2025-2030) provide concrete next steps.
\end{enumerate}

\subsection*{Key Principles}
\begin{itemize}
    \item \textbf{Evidence-Based Prioritization:} All priorities must pass the 4-Filter Framework (Regulatory Urgency, Foundational Dependency, Sovereignty Cliff, Market Feasibility).
    
    \item \textbf{Sovereignty-First Design:} Investments should systematically reduce foreign dependencies (measured via SSL metrics) while building European industrial capacity.
    
    \item \textbf{Compliance-by-Design:} All R\&I outputs must align with or anticipate NIS2, CRA, AI Act, and Data Act requirements to ensure market relevance.
    
    \item \textbf{Federated Interoperability:} Solutions must prioritize standardization and open architectures (e.g., Simpl middleware, open-source RISC-V) over proprietary, siloed approaches.
    
    \item \textbf{Measurable Progression:} Avoid subjective TRL assessments; instead, use objective ODP (market integration) and SSL (sovereignty) metrics with clear phase transition criteria.
\end{itemize}


% \section{Structure and main components of the Roadmap}
\section{Background \& Driving Factors}
% \noindent\textbf{Topic:} Cybersecurity for the Cognitive Computing Continuum\\
\textbf{Destinations:}
\begin{itemize}
    \item \textbf{Primary:} 4. A secure, sovereign European computing continuum infrastructure.
    \item \textbf{Secondary:} 6. European leadership in emerging disruptive computing paradigms.
\end{itemize}


\begin{itemize}
    \item \textbf{Decentralization Risk:} The paradigm shift from the centralized Cloud to a distributed Cloud-Edge-IoT continuum radically expands the attack surface. This creates environments that are often ``too many, too diverse, and too dispersed''  to secure with traditional perimeter defenses.
    \item \textbf{Physical Exposure:} Edge nodes and far-edge IoT devices frequently operate in the field. Unlike hyperscaler data centers, they lack strict physical and environmental security controls, making them highly vulnerable to tampering.
    \item \textbf{Regulatory Pressure:} There is an immediate, complex need to operationalize the NIS2 Directive \cite{nis2}, the Cyber Resilience Act (CRA) \cite{cra}, and the AI Act \cite{aiact} across a highly fragmented technological ecosystem to ensure compliance and market access.
    \item \textbf{Digital Sovereignty:} Europe must urgently reduce its systemic reliance on non-EU technology providers. Doing so will mitigate extra-territorial data access risks and avoid geopolitical vendor lock-in, ensuring independent control over critical infrastructure.
\end{itemize}

\subsection{Europe in Relation to the State of the Art}
\textbf{Comparison vs. Hyperscalers:}
\begin{itemize}
    \item \textbf{EU Strengths:} Europe leads globally in Trust, Privacy-by-Design, proactive regulation (GDPR, Data Act, AI Act), and rigorous, standardized certification frameworks (e.g., EUCS \cite{eucs}).
    \item \textbf{EU Gaps:} The European market suffers from a lack of unified, hyper-scale security orchestration. Furthermore, there is a critical dependency on non-EU hardware (secure silicon) and highly fragmented threat intelligence sharing among member states.
\end{itemize}
\textbf{Strategic Goal:} To definitively transition the European computing continuum from a state of ``reactive, siloed security''  to an advanced architecture defined by ``autonomous, federated, and certified Security-by-Design.''

